\section{第5章：静电场（5.5练习题）}
\sourceproblem{选择1}{80}{86}
\problemtext 两固定点电荷附近引入第三点电荷，判断原两者之间的相互作用力。
\workings 真空中点电荷之间的成对库仑力只取决于该对电荷及间距，故不变，选C。
\crosscheck 总合力一般变化；这里不考虑导体极化或原电荷位置改变。
\sourceproblem{选择2}{80}{86}
\problemtext 边长 $a$ 的立方体中心点电荷 $Q$，求顶点场强大小。
\workings 中心至顶点距离 $\sqrt3a/2$，故 $E=|Q|/(3\pi\varepsilon_0a^2)$，选C（按正电量约定）。
\crosscheck 场强大小非负，若 $Q<0$ 则方向向中心。
\sourceproblem{选择3}{80}{86}
\problemtext 原点正电荷 $Q$ 在 $(1,0)$ 产生场，求另一负电荷 $-2Q$ 的放置区域，使该点净场为零。
\workings 原场沿正x，负电荷必须在P左侧且共线，距离P为 $\sqrt2$（与坐标相同长度单位），故 $x=1-\sqrt2<0$，选C。
\crosscheck 仅匹配场强大小不足够，还须方向相反。
\sourceproblem{选择4}{80}{86}
\problemtext 半球开口平面与均匀电场平行，求曲面通量。
\workings 补底面后总通量为零，底面通量为零，故曲面通量也为零，选D。
\crosscheck 曲面各处场与法向夹角不同，不能直接用曲面积乘E。
\sourceproblem{选择5}{81}{87}
\problemtext 半球底面与均匀场夹30度，图中场有朝半球外侧的分量，求曲面外向通量。
\workings 底面外法向与场的夹角为120度，$\Phi_{\rm base}=-E\pi R^2/2$。闭合面净通量零，故 $\Phi_S=E\pi R^2/2$，选A。
\crosscheck 图中场箭头向右上，曲面净通量为正；若反转场方向则变号。
\sourceproblem{选择6}{81}{87}
\problemtext 立方体一顶点有点电荷 $q$，求图示不含电荷的对侧面通量。
\workings 拼合8个相同立方体使电荷在大立方体中心，每个大面通量 $q/(6\varepsilon_0)$，四个等价小面各得 $q/(24\varepsilon_0)$，选A。
\crosscheck 不直接对电荷位于边界的原立方体套内部电荷公式；以对称拼合规避奇点。
\sourceproblem{选择7}{81}{87}
\problemtext 两正点电荷相距 $2a$，以左电荷为中心作半径 $a$ 的高斯球，比较左右等面积小面元通量及总通量。
\workings 图中 $S_1$ 位于左端、$S_2$ 位于右端。中心电荷对两者贡献相等；右侧外电荷在左端的场沿外法向，在右端的场反外法向，因此 $\Phi_1>\Phi_2$；总通量 $q/\varepsilon_0$，选A。
\crosscheck 外部电荷不贡献闭合面净通量，却可改变局部通量。
\sourceproblem{选择8}{81}{87}
\problemtext 高斯面内电荷代数和为零，判断能确定的量。
\workings 高斯定理给 $\Phi_S=0$，仅此必然成立，选C。
\crosscheck 匀强场中任一闭面即为净通量零、局部场不为零的反例。
\sourceproblem{选择9}{81}{87}
\problemtext 判断关于高斯定理的四个推论。
\workings $\Phi=Q_{\rm enc}/\varepsilon_0$，若净内电荷非零则总通量非零，选D。净电荷为零不意味着内部无电荷；外场也可使面上各点场非零。
\crosscheck 对一般电场只由通量不能恢复逐点场强。
\sourceproblem{选择10}{81}{87}
\problemtext 无限长均匀带电圆柱面半径 $R$，辨认场强径向曲线（选项图在PDF88）。
\workings 轴对称高斯柱给出 $E=0$（$r<R$），$E=|\lambda|/(2\pi\varepsilon_0r)$（$r>R$），故在 $R$ 处跳变、外部反比衰减，选B。
\crosscheck 理想面电荷处场本身有两侧极限，不能当作连续体电荷。
\sourceproblem{选择11}{82}{88}
\problemtext 无限均匀正带电平面位于 $x=0$，辨认有向场强曲线。
\workings $E_x=(\sigma/2\varepsilon_0)\operatorname{sgn}(x)$，负半轴为负常量，正半轴为正同值常量，选C。
\crosscheck 平面处有跳变；电势连续。
\sourceproblem{选择12}{82}{88}
\problemtext 闭面包围点电荷Q，另一个点电荷由无穷远移到面外，判断总通量及面上场强。
\workings 净包围电荷不变，总通量仍为 $Q/\varepsilon_0$；外电荷场与原场叠加，一般使面上场强矢量改变，选D。
\crosscheck 外点电荷非零且位于有限距离时，其在面上任一点的场矢量均非零。
\sourceproblem{选择13}{82}{88}
\problemtext 图示非匀强场的电场线穿过球面，已知一小面通量 $\Delta\Phi_e$，求其余球面通量。
\workings 按图示球内无场线起止、无净电荷的教材意图，闭面总通量为零，故其余为 $-\Delta\Phi_e$，选A。
\crosscheck 若图仅是部分场线示意而并未给球内净电荷为零，则一般式为 $Q_{\rm enc}/\varepsilon_0-\Delta\Phi_e$，不能仅凭稀疏绘图证明无源。
\sourceproblem{选择14}{82}{88}
\problemtext 均匀带电球壳内总场为零，判断单个面元非零电荷在球内产生的场（选项续PDF89）。
\workings 在面元近似点电荷下，$\mathrm d\bm E=\mathrm dq\,\bm e_r/(4\pi\varepsilon_0r^2)$，球内任一点到面元距离有限且非零，因此不为零，选C。
\crosscheck 总场为零来自所有面元矢量叠加；假定 $\mathrm dq\ne0$。
\sourceproblem{选择15}{83}{89}
\problemtext 高斯面内净电荷为零时，判断面上场强。
\workings 总通量为零而场强不必处处为零，选C。反之总通量非零也不能保证面上每点的场强均非零。
\crosscheck 外部电场或内部等量异号电荷均可产生局部非零场。
\sourceproblem{选择16}{83}{89}
\problemtext 球面中心有点电荷，判断哪些操作改变闭面通量。
\workings 只有引入另一非零电荷到面内会改变净包围电荷，因此选B。
\crosscheck 移动内部点电荷或缩小仍包围它的球面不改变总通量。
\sourceproblem{选择17}{83}{89}
\problemtext 球对称场的图示为内区随半径线性上升、外区按平方反比下降，辨认电荷分布。
\workings 均匀体密度 $\rho$ 的实心球内 $E=|\rho|r/(3\varepsilon_0)$，外区 $E=|\rho|R^3/(3\varepsilon_0r^2)$，选B。
\crosscheck $\rho=Ar$ 对应内部二次函数，$\rho=A/r$ 对应内部常量（除中心奇点）；均不符合图示。
\sourceproblem{选择18}{83}{89}
\problemtext 正方形四顶点有等量电荷，中心场强和电势均为零，选符号排列。
\workings 取无穷远零势，四顶点至中心等距，电势为零要求两正两负。相对顶点同号时两对场分别抵消，故a、c为正、b、d为负的选项C正确。
\crosscheck 正负全部互换也成立；相邻顶点同号的排列净场非零。
\sourceproblem{选择19}{83}{89}
\problemtext 点电荷 $+q$ 至P距离 $a$，M距离 $2a$，以P为零势求M电势。
\workings $V_M=k_eq(1/(2a)-1/a)=-q/(8\pi\varepsilon_0a)$，选D。
\crosscheck P和M均在电荷右侧，远处电势低于近处。
\sourceproblem{选择20}{83}{89}
\problemtext 边长0.3 m的等边三角形两顶点分别带正负 $10^{-8}$ C，求第三顶点场强和电势（选项续PDF90）。
\workings 单个场强为 $9\times10^9(10^{-8})/0.3^2=1000$ V/m。两场夹120度，矢量和大小仍1000 V/m，方向平行正负电荷连线、由正指向负。电势等量反号抵消，$V=0$（无穷远零势），选B。
\crosscheck 电势是标量代数和，场强是矢量和。
